Duem a terme una experiència de ressonància en un tub amb aigua que consisteix a submergir un tub obert pels dos extrems en un recipient que conté aigua, tal com es mostra en la figura de la dreta. Damunt de l'extrem superior del tub fem vibrar un diapasó, que emet un so de freqüència 442~Hz. La longitud de la columna d'aire ($L$) s'ajusta elevant el tub fora de l'aigua fins a trobar un punt on es produeix la ressonància (se sent una nota intensa). Comencem l'experiència amb tot el tub submergit i observem la primera ressonància quan $L = 19,3$~cm i la segona ressonància quan $L = 58,0$~cm.

\figura[0.18]{fig1.png}

\begin{apartat}{1,25}
Dibuixeu la forma de l'ona ressonant per a la primera i segona ressonàncies. Indiqueu en tots dos casos de quin harmònic es tracta i identifiqueu els ventres i els nodes. Justifiqueu per què en tots dos casos la longitud d'ona no varia. Determineu la longitud del tub que ha de quedar per sobre de l'aigua quan ressona el cinquè harmònic.
\end{apartat}

\begin{apartat}{1,25}
A partir dels resultats de l'experiència, determineu la velocitat del so a l'aire. Si substituïm l'aigua per glicerina, variarà aquest resultat? Raoneu la resposta.
\end{apartat}

\dades{%
  La velocitat del so a l'aigua és de $1\,493\ \mathrm{m\,s^{-1}}$.\\
  La velocitat del so a la glicerina és de $1\,904\ \mathrm{m\,s^{-1}}$.}
